\chapter{Resultados e Discussão}
\label{cap:10}

Este capítulo apresenta os resultados obtidos com base nos requisitos funcionais e não funcionais definidos no Capítulo~\ref{cap:05} e discute, a partir do que foi verificado, em que medida a plataforma atende ao problema enunciado no Capítulo~\ref{cap:01}. A Seção~\ref{sec:plataformaOperacao} descreve a plataforma em operação e delimita o que esse fato sustenta; a Seção~\ref{sec:validacaoRF} apresenta a validação dos requisitos funcionais; a Seção~\ref{sec:validacaoRNF} trata da validação dos requisitos não funcionais, com os limiares e o alcance de cada verificação; a Seção~\ref{sec:validacaoUso} relata o que se observou no uso da plataforma publicada; e a Seção~\ref{sec:discussao} retoma o problema do Capítulo~\ref{cap:01} à luz desses resultados.

\section{A Plataforma em Operação}
\label{sec:plataformaOperacao}

A plataforma encontra-se em operação pública desde 3 de outubro de 2026, data em que a aplicação cliente e a API foram publicadas no ambiente descrito no Capítulo~\ref{cap:09}, com a aplicação cliente na Vercel, a API em uma máquina virtual da Oracle Cloud, atrás do Caddy, o banco de dados no Supabase e as imagens no Cloudflare R2. O envio de e-mail pelo servidor SMTP do Gmail está ativo em produção, e as mensagens enviadas pela plataforma chegam ao destinatário, o que completa os fluxos de acesso que dependem de um canal externo à própria aplicação.

Convém delimitar desde já o que esse fato sustenta e o que não sustenta. Trata-se de operação em caráter de validação acadêmica, e não de exploração do produto, de modo que os registros apresentados adiante caracterizam o funcionamento da plataforma e não constituem indicador de adoção. A distinção é necessária porque a existência de um endereço público poderia sugerir uma base de usuários que o trabalho não possui e não afirma possuir. As Categorias, Coleções, Itens e \textit{layouts} existentes na base foram criados, em sua maior parte, pela própria autora, no curso do desenvolvimento. O uso por outras pessoas restringe-se a duas pessoas do convívio da autora, sem vínculo com o desenvolvimento, cujas observações são relatadas na Seção~\ref{sec:validacaoUso} e cujo alcance é ali delimitado.

Dois conjuntos de registros sustentam a validação que se segue. O primeiro é a execução dos testes automatizados descrita no Capítulo~\ref{cap:08}, na qual, em 7 de outubro de 2026, os 56 conjuntos de testes de unidade do \textit{backend}, com 407 testes no total, foram aprovados. O segundo é uma medição do tempo de resposta da API, realizada no mesmo dia e sobre a mesma versão do código, com um acervo sintético criado para esse fim, composto de uma Categoria, duas Coleções, cinquenta Itens e um \textit{layout} de quinze Elementos. Cada uma das oito rotas medidas recebeu cinquenta requisições, precedidas de cinco de aquecimento e espaçadas de 700~ms, e nenhuma das 441 requisições terminou em erro.
% Medição de 07/10/2026 no desktop Ryzen 5 5500; dados brutos em scripts/medicoes/resultados/local-20261007T235528 (backend b3aaa96).
A API foi executada em modo de produção no computador descrito no Quadro~\ref{quad:equipamento}, com processador AMD Ryzen 5 5500 e 16~GB de memória, conectada ao banco de dados do Supabase pela internet; cada tempo medido inclui, portanto, a ida e a volta até o banco. Concluída a medição, a conta de teste e todo o acervo sintético foram excluídos, de modo que nenhum dado da medição permanece na base.

O que esses registros estabelecem, em conjunto, é modesto e preciso. Eles mostram que a plataforma está publicada e acessível, os fluxos que dependem de e-mail funcionam em produção, as regras de negócio do \textit{backend} passaram pelo conjunto de testes de unidade e a API responde, sobre um acervo de tamanho conhecido, em tempos que podem ser confrontados com as métricas do Capítulo~\ref{cap:05}. A verificação de cada requisito funcional é objeto da Seção~\ref{sec:validacaoRF}.

\section{Validação dos Requisitos Funcionais}
\label{sec:validacaoRF}

Os requisitos funcionais foram verificados por três instrumentos complementares. O primeiro são os testes de unidade descritos no Capítulo~\ref{cap:08}, cuja distribuição por requisito consta do Quadro~\ref{quad:resumoTestes} e cujos casos principais, CT01 a CT08, especificam as regras de maior risco. O segundo é a conferência da implementação no código da aplicação cliente e da API, que estabelece por qual tela e por qual rota cada requisito se realiza. O terceiro é o uso da plataforma publicada, que exercita em conjunto a interface, a API e os serviços externos. A divisão de trabalho entre eles decorre do alcance de cada um, já que os testes de unidade verificam as regras de negócio do \textit{backend} com o banco de dados simulado, e o que depende da interface ou de um serviço externo, como o envio de e-mail, só se demonstra pela inspeção e pelo uso. As subseções seguem os grupos de requisitos definidos no Capítulo~\ref{cap:05}.

\subsection{Acesso ao Sistema}

O cadastro (RF-01) gera o código de usuário único e entrega à conta nova as seis paletas predefinidas e os \textit{layouts} padrão; o CT06 verifica que a paleta em uso nasce clara ou escura conforme o tema do sistema operacional informado no cadastro. A conta só é acessada depois da confirmação do e-mail por um código de seis dígitos, válido por quinze minutos e limitado a cinco tentativas, do qual o banco guarda apenas o \textit{hash}. A autenticação (RF-02) registra uma sessão no banco a cada \textit{login} e a confere a cada requisição, e o CT01 verifica que o \textit{logout} revoga apenas a sessão em uso e que um \textit{token} de sessão revogada recebe a resposta 401. A recuperação de senha (RF-04) é coberta pelo CT02, que verifica a mensagem idêntica para e-mails cadastrados ou não, o \textit{token} aleatório guardado como \textit{hash}, com prazo de validade e uso único, e a revogação de todas as sessões ao fim da redefinição; o envio de e-mail, do qual o \textit{link} depende, funciona em produção, conforme a Seção~\ref{sec:plataformaOperacao}. A alteração de senha (RF-03) é realizada pela API, que exige a senha atual, recusa uma nova senha igual à anterior e revoga as demais sessões na mesma transação, comportamento coberto pelos testes do módulo \texttt{user}. A edição do perfil (RF-05) ocorre em uma janela sobreposta ao próprio perfil, com a foto enquadrada em círculo e a capa enquadrada na proporção da faixa, ambas com o original guardado para reenquadramento.

\subsection{Coleção}

A organização em Categorias, Coleções e Itens (RF-06, RF-07, RF-09 e RF-13) foi verificada pelos testes dos módulos \texttt{category}, \texttt{collection} e \texttt{item} e pelo uso. O painel de Categorias exibe as contagens de Coleções e de Itens, alterna entre grade e lista e oferece as quatro ordenações previstas, e cada nível impede nomes repetidos no seu escopo. O CT08 verifica que a imagem enviada é reconhecida pela assinatura do arquivo, e não pela extensão declarada, de modo que um arquivo de texto renomeado é recusado. As \textit{tags} são gravadas em forma canônica e formam um vocabulário único, e a tela de descoberta reúne as Coleções de outros usuários que tenham todas as \textit{tags} escolhidas, respeitando a privacidade de cada uma. O controle de privacidade (RF-08) é o requisito com verificação mais extensa, pois o CT04 verifica o nível de acesso do dono, do amigo e de terceiros a cada privacidade, com a resposta 404 para a Coleção que o visitante não pode ver, e o CT03 verifica que restringir a privacidade encerra, na mesma transação, o acompanhamento de quem perdeu o acesso.

A movimentação de Itens (RF-10) é concluída em três comandos pelo menu de ações do Item, em que o usuário escolhe a ação, escolhe o destino e confirma. Durante trinta segundos, um botão permite desfazê-la, e os testes do módulo \texttt{item} verificam que o Item com ficha própria a conserva ao mudar de Coleção. A ficha em modo leitura (RF-11) é aberta sobre a listagem, com setas que percorrem os Itens na ordem exibida, indicação visual dos Elementos vazios e nenhuma opção de edição para o visitante. Por ser um requisito de apresentação, sem regra de negócio própria no \textit{backend}, o RF-11 não tem teste de unidade associado no Quadro~\ref{quad:resumoTestes} e foi verificado pela inspeção e pelo uso. A listagem de Itens (RF-12) oferece busca pelo nome e pelo subtítulo, as quatro ordenações por nome e por data e a opção de esconder o texto sobre as capas, e o cartão exibe a média normalizada dos Elementos de Classificação da ficha.

\subsection{Editor de \textit{Layout}}

O editor (RF-14) oferece os seis tipos de Elemento previstos e posiciona cada um em unidades da folha, cujo tamanho é independente dos Elementos. O editor e a ficha em modo leitura desenham o \textit{layout} com o mesmo componente, de modo que a prévia é, por construção, a própria ficha final, e não uma aproximação dela. A duplicação cria uma cópia independente, com o sufixo ``(cópia)'', e a exclusão de um \textit{layout} em uso é precedida da desvinculação, que entrega antes a cada Item uma cópia própria do \textit{layout}, sem alterar nenhuma ficha. A associação a Coleções (RF-15) admite a troca com confirmação e a mescla das fichas próprias no novo \textit{layout}, cuja regra de correspondência (mesma origem, mesmo tipo e nome e, por fim, mesmo tipo na ordem de leitura) é verificada pelos testes da mescla no módulo \texttt{common}. O compartilhamento (RF-16) importa um \textit{layout} a partir do seu identificador e cria uma cópia independente marcada como importada, e os testes do módulo \texttt{layout} verificam a recusa de importar um \textit{layout} da própria biblioteca. A personalização de um Item (RF-17) é verificada pelo CT05, segundo o qual a primeira edição cria a cópia própria da ficha, com ou sem \textit{layout} na Coleção, e uma edição sobre a ficha ainda não copiada é recusada com a resposta 409.

\subsection{Aparência}

As paletas (RF-18) são aplicadas a toda a interface no momento da escolha e gravadas no servidor, o que as mantém entre sessões e dispositivos. O CT06 verifica a troca da paleta em uso como operação atômica e a recusa de excluir a paleta em uso, e os testes do módulo \texttt{appearance-preset} verificam que cada paleta tem a cor primária e as doze subcores preenchidas. O editor de paletas calcula a razão de contraste de cada par de cor de texto e de superfície e avisa quando ela fica abaixo de 4,5:1, sem impedir a escolha. A tipografia (RF-19) integra a paleta e oferece dezenove famílias de fonte, com os tamanhos, o espaçamento entre linhas e o peso validados nas faixas do requisito pelos testes do DTO correspondente.

\subsection{Interação Social}

A amizade (RF-20) é solicitada a partir do código de usuário, e o CT07 verifica que só uma solicitação pendente pode ser aceita, que a notificação ao solicitante é criada na mesma transação do aceite e que dois aceites simultâneos resultam em um único vínculo, com a resposta 409 para o segundo; a recusa, por sua vez, não gera notificação. O acompanhamento de Coleções (RF-21) obedece às mesmas regras de visibilidade da privacidade, verificadas pelo CT03 e pelos testes do módulo \texttt{social}, e a tela de Coleções seguidas reúne, para cada uma, o dono, a Categoria e a quantidade de Itens. A visita a outros usuários (RF-23) percorre o perfil, a Categoria e a página da Coleção em modo somente leitura; o CT04 verifica a resposta 404 para o que o visitante não pode ver, e o perfil visitado omite as Categorias sem nenhuma Coleção visível. O \textit{feed} de atualizações (RF-22) é fornecido pela API, que devolve os eventos das Coleções acompanhadas, do mais recente ao mais antigo, com paginação por cursor, excluindo as Coleções privadas e preservando o nome do Item já removido, comportamento coberto pelos testes do serviço de \textit{feed}.

\subsection{Notificação}

As notificações (RF-24) ficam no sino do cabeçalho, presente em todas as telas da área logada, com a quantidade de não lidas, consultada a cada sessenta segundos e sempre que o usuário retorna à janela. Abrir uma notificação a marca como lida e leva ao que a originou, seja a Coleção atualizada, seja o perfil de quem pediu ou aceitou a amizade. Os testes do módulo \texttt{notification} verificam a marcação individual e coletiva, a limpeza da área sem exclusão dos registros e a ocultação das notificações de Coleções que passaram a ser privadas. O canal de e-mail, usado pela confirmação de cadastro e pela recuperação de senha, funciona em produção, conforme a Seção~\ref{sec:plataformaOperacao}.

O Quadro~\ref{quad:sinteseRF} reúne o instrumento, o resultado e a situação de cada requisito funcional.

\begin{quadro}[h!]
\caption{Síntese da validação dos requisitos funcionais}
\label{quad:sinteseRF}
\footnotesize
\begin{tabular}{|p{1.3cm}|p{4.6cm}|p{5.6cm}|p{2.2cm}|}
\hline
\textbf{Req.} & \textbf{Instrumento/Evidência} & \textbf{Resultado} & \textbf{Situação} \\ \hline
RF-01 & CT06; módulo \texttt{user}; inspeção & Código de usuário, paletas e \textit{layouts} iniciais; confirmação do e-mail por código & Atendido \\ \hline
RF-02 & CT01; módulo \texttt{auth} & Sessão revogada recebe 401 & Atendido \\ \hline
RF-03 & Módulo \texttt{user} & Senha atual exigida; demais sessões revogadas & Implementado na API \\ \hline
RF-04 & CT02; uso em produção & \textit{Token} de uso único; \textit{link} entregue por e-mail & Atendido \\ \hline
RF-05 & Módulo \texttt{user}; inspeção & Foto e capa enquadradas, com original guardado & Atendido \\ \hline
RF-06 & Módulo \texttt{category}; uso & Categoria com capa, ícone e tema & Atendido \\ \hline
RF-07 & Módulo \texttt{collection}; uso & Coleções com \textit{tags} e descoberta por \textit{tags} & Atendido \\ \hline
RF-08 & CT03; CT04 & Acesso por privacidade; acompanhamentos encerrados & Atendido \\ \hline
RF-09 & CT08; módulo \texttt{item} & Imagem reconhecida pela assinatura & Atendido \\ \hline
RF-10 & Módulo \texttt{item}; uso & Três comandos; desfazer por 30~s; ficha própria mantida & Atendido \\ \hline
RF-11 & Inspeção; uso & Ficha sobreposta, somente leitura ao visitante & Atendido \\ \hline
RF-12 & Módulo \texttt{item}; uso & Busca, ordenação e média das Classificações & Atendido \\ \hline
RF-13 & Módulo \texttt{category}; uso & Contagens, grade e lista, ordenações & Atendido \\ \hline
RF-14 & Módulo \texttt{layout}; inspeção & Seis tipos de Elemento; prévia igual à ficha & Atendido \\ \hline
RF-15 & Testes da mescla; uso & Troca com confirmação e mescla das fichas & Atendido \\ \hline
RF-16 & Módulo \texttt{layout} & Cópia independente importada pelo identificador & Atendido \\ \hline
RF-17 & CT05 & Cópia própria na primeira edição & Atendido \\ \hline
RF-18 & CT06; módulo \texttt{appearance-preset} & Troca atômica; aviso de contraste & Atendido \\ \hline
RF-19 & Testes do DTO de paleta & Faixas tipográficas validadas & Atendido \\ \hline
RF-20 & CT07 & Aceite com notificação na mesma transação & Atendido \\ \hline
RF-21 & CT03; módulo \texttt{social} & Acompanhamento conforme a visibilidade & Atendido \\ \hline
RF-22 & Testes do serviço de \textit{feed} & Eventos paginados, sem Coleções privadas & Implementado na API \\ \hline
RF-23 & CT04; uso & Visita somente leitura; 404 ao que não é visível & Atendido \\ \hline
RF-24 & Módulo \texttt{notification}; uso & Sino com contador; ocultação por privacidade & Atendido \\ \hline
\end{tabular}
\fonte{Elaborado pelo autor}
\end{quadro}

Os vinte e quatro requisitos funcionais têm, portanto, implementação verificada, sendo vinte e dois de ponta a ponta, da interface ao banco, e dois, o RF-03 e o RF-22, por rotas da API cobertas por testes de unidade. A seção seguinte examina se a plataforma apresenta as propriedades de desempenho, segurança e acessibilidade que os requisitos não funcionais enunciam.

\section{Validação dos Requisitos Não Funcionais}
\label{sec:validacaoRNF}

Os requisitos não funcionais do Capítulo~\ref{cap:05} não são o eixo deste trabalho, cuja contribuição está nas funcionalidades verificadas na seção anterior, mas cada um deles merece um juízo explícito sobre o seu atendimento. A situação de cada requisito assume um de quatro valores: \textbf{atendido}; \textbf{atendido com ressalva}, quando a métrica foi cumprida, mas alguma restrição do requisito comprovadamente não foi; \textbf{não verificado}, quando faltou o instrumento ou a medição; e \textbf{não atendido}, quando o resultado contraria a métrica. A situação é determinada pela métrica, e uma restrição não cumprida impede que o requisito seja considerado plenamente atendido, ainda que a métrica o seja.

\subsection{Método das Medições}

As medições foram feitas com \textit{scripts} mantidos no repositório da API, na pasta \texttt{scripts/medicoes}, que gravam cada requisição individualmente (horário, rota, código de resposta, tempo e tamanho), a descrição do ambiente e os resumos estatísticos, de modo que todo número citado nesta seção pode ser conferido e a medição, repetida. Os \textit{scripts} criam, pela própria API, a conta de teste e o acervo sintético descritos na Seção~\ref{sec:plataformaOperacao}, cujo \textit{layout} de quinze Elementos corresponde à quantidade citada no RNF-07, e cinco Itens com ficha própria. A confirmação do e-mail da conta de teste passou pela mesma rota usada pela interface, com o código lido do registro da API, executada em um modo que não envia e-mail. Ao final, os \textit{scripts} excluem a conta e confirmam a exclusão.

O tempo de cada requisição foi medido do instante anterior ao envio até a leitura completa do corpo da resposta, uma requisição de cada vez, com 700 milissegundos entre elas. O intervalo mantém o ritmo abaixo do limite de 100 requisições por minuto que a API impõe a cada endereço de origem; pela mesma razão, não foi feito ensaio de carga, que, a partir de uma única máquina, mediria o limitador, e não a API. Os percentis foram calculados pelo método do posto mais próximo, sem interpolação, de modo que todo valor relatado é um tempo que de fato ocorreu; com 50 amostras, o percentil 99 coincide com o máximo. O Quadro~\ref{quad:latenciaApi} apresenta os resultados.

\FloatBarrier
\begin{quadro}[h!]
\caption{Tempo de resposta da API, em milissegundos (n = 50 por rota)}
\label{quad:latenciaApi}
\footnotesize
\begin{tabular}{|p{6.0cm}|p{1.3cm}|p{1.4cm}|p{1.3cm}|p{1.3cm}|p{1.4cm}|}
\hline
\textbf{Rota} & \textbf{RNF} & \textbf{Mínimo} & \textbf{p50} & \textbf{p95} & \textbf{Máximo} \\ \hline
Painel de Categorias (\texttt{GET /category}) & RNF-02 & 58,3 & 62,8 & 72,3 & 77,8 \\ \hline
Coleções da Categoria (\texttt{GET /collection}) & RNF-02 & 77,1 & 83,2 & 103,5 & 128,0 \\ \hline
Lista de 50 Itens (\texttt{GET /item}) & RNF-02 & 110,8 & 118,1 & 135,3 & 150,9 \\ \hline
Ficha própria com 15 Elementos (\texttt{GET /item/:id/layout}) & RNF-02 & 87,3 & 95,3 & 109,6 & 117,4 \\ \hline
Cadastro de Item (\texttt{POST /item}) & RNF-02 & 251,2 & 268,9 & 305,7 & 317,9 \\ \hline
Exclusão de Item (\texttt{DELETE /item/:id}) & RNF-02 & 143,4 & 155,8 & 169,7 & 194,7 \\ \hline
Mover Item entre Coleções (\texttt{PATCH /item/:id}) & RNF-06 & 278,1 & 296,1 & 326,4 & 337,6 \\ \hline
Gravação de Elemento no editor (\texttt{PATCH /layout/:id/elements/:id}) & RNF-07 & 90,6 & 99,4 & 107,6 & 118,7 \\ \hline
\end{tabular}
\fonte{Elaborada pelo autor}
\end{quadro}
\FloatBarrier

A medição não foi repetida no ambiente de produção descrito no Capítulo~\ref{cap:09}, de modo que os tempos acima caracterizam a API e o banco de dados, mas não a distância entre a máquina virtual e o Supabase nem a rede de quem usa a plataforma publicada.
% [SE MEDIR EM PRODUÇÃO: latencia-api.mjs com AMBIENTE=producao; acrescentar o quadro equivalente aqui]

\subsection{Requisitos Mensuráveis}

\textbf{RNF-02 (desempenho de carregamento).} O requisito fixa três limites: três segundos para o carregamento de uma tela principal, dois segundos para as respostas de leitura da API e três segundos para as operações de escrita. Para a tela, a medida adequada seria o maior elemento de conteúdo pintado (LCP), no percentil 75 dos carregamentos, convenção que acompanha a recomendação do web.dev, para o qual um LCP de até 2,5 segundos caracteriza um bom carregamento \cite{walton2019lcp}. Para a API, adota-se o percentil 95, que afirma o limite para a quase totalidade das requisições, e não apenas para a típica. O maior percentil 95 entre as rotas de leitura foi de 135,3 milissegundos (lista de 50 Itens), e o maior entre as de escrita, de 305,7 milissegundos (cadastro de Item), valores cerca de quinze e de dez vezes menores que os respectivos limites. O carregamento das telas, que depende do navegador, não foi medido. Independentemente dessa medição, duas restrições do requisito não são atendidas: as imagens não são otimizadas pelo servidor, e as listagens de Categorias, Coleções e Itens não são paginadas, sendo devolvidas por inteiro em uma única resposta. O requisito permanece \textbf{não verificado}, e, ainda que o carregamento das telas cumpra o limite, a situação máxima será a de atendido com ressalva.
% [SE MEDIR AS TELAS: LCP p75 de dez carregamentos da tela de Itens e da ficha, com scripts/medicoes/tela.js]

\textbf{RNF-06 (movimentação de Itens).} O requisito exige que a mudança de Coleção se reflita na interface em até 500 milissegundos e que a interface a aplique de forma otimista, antes da confirmação do servidor. A implementação não é otimista: a interface aguarda a resposta da API e, em seguida, recarrega a listagem, e só então o Item deixa de aparecer na Coleção de origem. O tempo percebido é, portanto, no mínimo a soma da gravação com a nova leitura da lista, à qual se acrescenta o desenho da tela, que não foi medido. Na medição, essa soma foi de 388,9 milissegundos com os mínimos e de 461,7 com os percentis 95, abaixo do limite, mas com folga de apenas 38,3 milissegundos no segundo caso. A margem é estreita o bastante para depender do ambiente: em uma medição preliminar, feita em 2 de outubro de 2026 em outro computador e com o mesmo banco, a soma dos mínimos foi de 568,3 milissegundos, acima do limite. A gravação de uma movimentação é a mais lenta das rotas medidas, o que se explica pelo número de operações executadas em sequência, cada uma atravessando a internet até o banco: a verificação da sessão, a confirmação de que o Item e a Coleção de destino pertencem ao usuário e a verificação de unicidade do nome, seguidas de uma transação que atualiza o Item, registra o evento e consulta quem acompanha a Coleção. Como a soma medida é apenas um limite inferior do tempo percebido, o requisito é considerado \textbf{não verificado}, e a restrição de aplicação otimista não é atendida, o que limita a situação máxima à de atendido com ressalva. A restrição de desfazer foi atendida de forma diferente da especificada: a interface oferece, por 30 segundos, um botão que devolve o Item à Coleção de origem, e não o atalho \texttt{Ctrl+Z}.

\textbf{RNF-07 (desempenho do editor).} O requisito limita a 200 milissegundos o tempo entre uma operação no editor e sua exibição. Esse valor coincide com o limite que o web.dev adota para uma boa responsividade na métrica INP, avaliada no percentil 75 das interações \cite{wagner2022inp}, que seria, por isso, a medida adequada. A gravação no servidor não compõe esse tempo, pois ocorre de forma assíncrona, depois de uma pausa de 400 milissegundos, sem bloquear a edição, o que atende à restrição de salvamento assíncrono; ainda assim, registra-se que ela levou 107,6 milissegundos no percentil 95. A responsividade das interações no editor não foi medida, e o requisito permanece \textbf{não verificado}.
% [SE MEDIR O EDITOR: p75, p98 e máximo de pelo menos 20 interações, com scripts/medicoes/editor.js]

\textbf{RNF-08 (disponibilidade).} A disponibilidade de 99\% só pode ser aferida por um serviço externo que consulte periodicamente a plataforma publicada, e esse acompanhamento não foi feito; o requisito permanece \textbf{não verificado}. A restrição de recuperação automática é atendida pela política de reinício do contêiner e pelos serviços habilitados na inicialização da máquina virtual, descritos no Capítulo~\ref{cap:09}. Duas restrições não são atendidas independentemente da medição: não há cópia local para uso sem conexão, e a interface não reenvia automaticamente operações feitas enquanto a conexão estava interrompida, com exceção das alterações do editor, cuja sincronização reenvia o que ficou pendente na tentativa seguinte.

\subsection{Síntese da Validação}

O Quadro~\ref{quad:sinteseRNF} reúne a situação de todos os requisitos não funcionais, inclusive os que não dependem de medição de tempo.

\FloatBarrier
\begin{quadro}[h!]
\caption{Síntese da validação dos requisitos não funcionais}
\label{quad:sinteseRNF}
\footnotesize
\begin{tabular}{|p{2.3cm}|p{3.3cm}|p{5.7cm}|p{2.2cm}|}
\hline
\textbf{Requisito} & \textbf{Instrumento} & \textbf{Resultado} & \textbf{Situação} \\ \hline
RNF-01 Usabilidade & Teste com usuários iniciantes (não realizado) & Nenhuma medição do tempo para criar uma Coleção e um Item & Não verificado \\ \hline
RNF-02 Carregamento & \textit{Script} de latência & Leitura: p95 máximo de 135,3\,ms; escrita: p95 máximo de 305,7\,ms. Telas não medidas. Sem paginação e sem otimização de imagens & Não verificado \\ \hline
RNF-03 Segurança de senhas e dados & Inspeção do código; testes automatizados (Capítulo~\ref{cap:08}) & bcrypt, mensagem genérica no \textit{login}, sessões revogáveis, HTTPS em produção e limite de cinco \textit{logins} por minuto. Registro de erros inesperados com o objeto completo da exceção & Atendido com ressalva \\ \hline
RNF-04 Compatibilidade & Teste em navegadores e larguras de tela (não realizado) & Interface escrita para telas a partir de 320\,px; matriz de navegadores não testada & Não verificado \\ \hline
RNF-05 Personalização & Inspeção do código e da interface & Variáveis CSS, alerta de contraste, aplicação imediata e persistência no perfil; combinações não testadas sistematicamente & Atendido com ressalva \\ \hline
RNF-06 Movimentação de Itens & \textit{Script} de latência; análise do fluxo da interface & Gravação mais leitura: 461,7\,ms no p95; desenho não medido; atualização não otimista; desfazer por 30\,s via botão & Não verificado \\ \hline
RNF-07 Editor & \textit{Script} de latência & Gravação assíncrona com p95 de 107,6\,ms; responsividade das interações não medida & Não verificado \\ \hline
RNF-08 Disponibilidade & Monitoramento externo (não realizado) & Reinício automático configurado; sem medição; sem cópia local para uso sem conexão & Não verificado \\ \hline
RNF-09 Acessibilidade & Inspeção do código e registros do projeto & Contraste das paletas padrão verificado e alerta de contraste; navegação do editor por teclado não implementada & Não atendido \\ \hline
RNF-10 Integridade & Inspeção do código; configuração do banco & Transações atômicas e confirmação antes de excluir; sem cópia de segurança automática no plano gratuito; sem registro de auditoria das exclusões & Não atendido \\ \hline
\end{tabular}
\fonte{Elaborada pelo autor}
\end{quadro}
\FloatBarrier

Três situações merecem comentário. O RNF-03 é considerado atendido com ressalva porque os mecanismos que a métrica exige existem e estão verificados, mas a restrição de não expor dados sensíveis em registros do sistema não é cumprida por inteiro: os erros inesperados são registrados com o objeto completo da exceção, que pode conter valores de campos, e a falha no envio de um código de confirmação registra o e-mail do destinatário. O limite de tentativas de \textit{login} por endereço, por sua vez, depende, atrás do Caddy, da configuração de confiança no intermediário descrita no Capítulo~\ref{cap:09}.
% [DEPENDE DE CORREÇÃO: app.set('trust proxy', 1) no main.ts; correcoes_sistema.md, itens 2 e 10]
O RNF-09 é considerado não atendido porque o requisito exige conformidade com o nível AA das WCAG e que todos os elementos interativos sejam acessíveis por teclado, enquanto a navegação do editor por teclado não foi implementada, de modo que posicionar um Elemento depende de um dispositivo apontador ou dos campos numéricos de posição; o critério de operação por teclado é de nível A, abaixo do AA exigido \cite{campbell2024wcag}. As demais partes da interface não foram auditadas. Por fim, o RNF-10 é considerado não atendido porque a sua métrica principal, a de cópias de segurança diárias com retenção de 30 dias, não é cumprida no plano gratuito do Supabase, que não faz cópia automática e deixa a exportação dos dados a cargo do responsável pelo projeto \cite{supabase2026backups}, e porque o registro em auditoria de toda exclusão também não é atendido: apenas as alterações de Itens em Coleções não privadas geram registro de evento.
% [VERIFICAR com a autora: se houver rotina de db dump na VM (cap. 9), rever o RNF-10]

O resultado de conjunto é coerente com a posição secundária desses requisitos no trabalho. Os tempos de resposta da API, que dependem apenas do servidor, ficaram com folga dentro dos limites; a movimentação de Itens, que depende também da interface, ficou dentro do limite por margem estreita e sem a atualização otimista exigida. Os requisitos que exigem avaliação com pessoas, como usabilidade e acessibilidade, ou acompanhamento da plataforma publicada, como a disponibilidade, permanecem em aberto, e são essas as lacunas que a avaliação deixa.

\section{Validação em Uso}
\label{sec:validacaoUso}

A validação em uso foi conduzida por observação, sem entrevista nem questionário, e registra o que ocorreu quando a plataforma publicada foi usada pela autora e por duas pessoas do seu convívio que não participaram do desenvolvimento e que, por isso, não conheciam a plataforma de antemão.

Desse uso resultou uma observação contrária ao previsto. Uma das usuárias, diante da versão publicada, pediu um guia de primeiros passos, e o pedido revelou que ela esperava uma ferramenta de listas, próxima de uma planilha. O modelo da plataforma, em que um \textit{layout} reutilizável é associado a uma Coleção e dá forma às fichas dos seus Itens, não se mostrou evidente sem orientação. A observação tem peso particular porque a planilha é justamente uma das ferramentas que o Capítulo~\ref{cap:01} caracteriza como rígidas, por reduzirem cada registro a campos fixos. O primeiro contato com a plataforma reproduziu, portanto, a expectativa formada por essas ferramentas, e não a proposta que pretende substituí-las.

Convém registrar a ordem dos fatos. A decisão de oferecer um guia de primeiros passos antecede esse pedido, de modo que a observação converge com ela sem tê-la motivado. O guia compreende uma apresentação no primeiro acesso, com um diagrama que relaciona o \textit{layout}, a Coleção e as fichas, uma lista de primeiros passos acompanhada de dicas em cada tela e uma página que explica o funcionamento da plataforma. Sobre a eficácia desse guia nada se afirma, pois ele não foi observado em uso por quem fez o pedido nem por outra pessoa.

O alcance destas observações é restrito, já que envolvem três pessoas, nenhuma independente do trabalho, e um único relato que vai além do próprio uso. Nada se afirma quanto à facilidade de uso em geral, à preferência pela plataforma ou à sua adoção, e a Seção~\ref{sec:discussao} retoma esses limites.

\section{Discussão e Implicações}
\label{sec:discussao}

O Capítulo~\ref{cap:01} enunciou o problema que motivou este trabalho. As ferramentas usualmente empregadas para registrar coleções pessoais tratam cada registro como um conjunto de campos fixos com aparência padronizada, e falta uma plataforma que reúna a organização de coleções de qualquer tipo, a liberdade de definir a estrutura e a aparência de cada registro e a possibilidade de compartilhar essas coleções com outras pessoas. Esta seção retoma esse problema à luz do que as seções anteriores apuraram.

A resposta é afirmativa quanto ao que foi verificado e restrita quanto ao que se pode generalizar. Cada uma das três partes do problema tem correspondência na plataforma em operação.

\begin{itemize}
    \item \textbf{Organização de coleções de qualquer tipo:} realiza-se na hierarquia de Categorias, Coleções e Itens, cujo tema é texto livre definido pelo usuário.
    \item \textbf{Liberdade de estrutura e de aparência:} realiza-se no editor, em que o próprio usuário define quais Elementos compõem a ficha e como eles se dispõem, e na personalização global por paletas e tipografia.
    \item \textbf{Compartilhamento:} realiza-se pela visita a perfis, pelo acompanhamento de Coleções e pelas notificações, sob uma privacidade definida em cada Coleção e verificada pelos testes de maior extensão do conjunto, conforme a Seção~\ref{sec:validacaoRF}.
\end{itemize}

Convém, contudo, separar o que esses resultados demonstram do que não demonstram. Eles demonstram que a proposta é construível e operável, porque os requisitos funcionais têm implementação verificada e a plataforma funciona publicada, com os serviços externos de que depende. Não demonstram que a plataforma seja preferível às planilhas, aos aplicativos de listas ou aos cadernos mencionados no Capítulo~\ref{cap:01}, porque nenhuma comparação de uso foi conduzida, nem que seja de uso fácil para o público em geral, porque as observações da Seção~\ref{sec:validacaoUso} envolvem três pessoas próximas do trabalho. A contribuição situa-se, por isso, no plano da demonstração de viabilidade, e não no da evidência de superioridade.

A observação de uso permite, ainda assim, uma implicação que extrapola esta plataforma. A flexibilidade que responde ao problema do Capítulo~\ref{cap:01} tem um custo de entendimento, pois um registro cuja estrutura é definida pelo próprio usuário exige que ele compreenda antes a relação entre o molde e as fichas, ao passo que a planilha, por ter estrutura fixa, dispensa essa compreensão. Para uma ferramenta desse tipo, a apresentação do modelo deixa de ser material acessório e passa a integrar o produto, o que corresponde à decisão de oferecer um guia de primeiros passos. Cabe registrar, como conjectura, que a expectativa de encontrar uma planilha seja comum entre usuários novos, por ser essa a referência mais difundida para registrar listas. A conjectura não pode ser verificada a partir de um único relato, e sua verificação exigiria observar usuários independentes, com e sem o guia.

As decisões de projeto que sustentam o editor também merecem ser qualificadas pelo que custam. O editor e a ficha em modo leitura desenham o \textit{layout} com o mesmo componente, o que torna a prévia fiel por construção, e não por verificação. Em contrapartida, toda a ficha precisa ser expressa em uma unidade relativa à largura da folha, resolvida no navegador por consultas de contêiner, recurso recente das folhas de estilo, e a fidelidade da prévia passa a depender desse suporte. A ficha própria de um Item, por sua vez, é uma cópia do \textit{layout} da Coleção, o que garante que personalizar um Item nunca altere os demais. O custo dessa escolha é que a cópia deixa de acompanhar as alterações do \textit{layout} de origem, e reaproximá-las exige a mescla descrita na Seção~\ref{sec:validacaoRF}, cuja correspondência entre Elementos é heurística e mantém na ficha o que não encontra lugar no novo \textit{layout}.

Observação análoga cabe à privacidade. A regra de visibilidade não reside em um ponto único do sistema, porque é imposta em cada rota que recebe o identificador de uma Coleção, de um Item ou de um perfil, e responde com 404 ao que o visitante não pode ver, de modo a não revelar sequer a existência do recurso. Na prática, a privacidade custa mais para manter do que para projetar, já que toda rota nova precisa aplicá-la, e esse custo cresce com a quantidade de rotas, e não com a quantidade de níveis de privacidade. Por fim, a publicação em serviços gratuitos, descrita no Capítulo~\ref{cap:09}, viabilizou a operação sem custo financeiro em troca de limites de recurso, uma vez que a API dispõe de 1~GB de memória e a recuperação dos dados depende de uma exportação feita pelo responsável pela plataforma. Para uma operação de validação, a troca é favorável, mas é uma troca, e não um ganho sem contrapartida.

Os tempos medidos permitem qualificar ainda uma escolha da interface. Ao aguardar a confirmação do servidor antes de retirar da Coleção de origem o Item movido, a interface não exibe um estado que o banco não tenha confirmado, o que dispensa reverter a tela quando a gravação falha. O custo é que o tempo percebido passa a somar as idas e voltas até o banco, e a Seção~\ref{sec:validacaoRNF} mostrou que essa soma cabe no limite do RNF-06 por margem estreita, que variou com o computador usado na medição. A atualização otimista inverteria a troca, pois o tempo do servidor deixaria de compor o tempo percebido, ao preço de desfazer a alteração na tela sempre que a gravação fosse recusada.
